\uuid{hUx2}
\titre{Système linéaire à un paramètre/Méthode du Pivot de Gauss}
\auteur{hanani ; mzari}
\organisation{exo7}
\licence{CC BY-NC-SA 4.0}
\datecreate{2020-01-20}
\niveau{L1}
\module{Algèbre}
\chapitre{Déterminant, système linéaire}
\sousChapitre{Système linéaire, rang}
\difficulte{3}
\theme{Systèmes d'équations linéaires, Système linéaire à un paramètre, Méthode du Pivot de Gauss}
\exo7qcmid{lille-2.systemes.014}

\contenu{
\begin{question}
\qclassification{200.03}
\qtags{motcle=Système linéaire à un paramètre/Méthode du Pivot de Gauss}
On considère le système d'équations, d'inconnue $(x,y,z)\in \Rr^3$ et de paramètre un réel  $m$ :
$$(\mathtt{S})
\left\{\begin{array}{rcc}
x+y+mz&=&1\\
x+my+z&=&1\\
mx+y+z&=&1.\\
\end{array}\right.$$
Quelles sont les assertions vraies ?
\begin{answers}
\bad{$(\mathtt{S})  \Leftrightarrow
\left\{\begin{array}{rcc}
x+y+mz&=&1\\
(m-1)y+(1-m)z&=&0\\
(1-m)z&=&1-m.\\
\end{array}\right.$}
\good{Si $m =1$,  $(\mathtt{S})$ admet une infinité de solutions.}
\good{Si $m=-2$, $(\mathtt{S})$ n'admet pas de solution.}
\bad{Si $m\neq 1$,  $(\mathtt{S})$ admet une unique solution.}
\end{answers}
\begin{explanations}
$$(\mathtt{S}) \Leftrightarrow
\left\{\begin{array}{rcc}
x+y+mz&=&1\\
(m-1)y+(1-m)z&=&0\\
(1-m)y+(1-m^2)z&=&1-m\\
\end{array}\right. \Leftrightarrow
\left\{\begin{array}{rcc}
x+y+mz&=&1\\
(m-1)y+(1-m)z&=&0\\
(1-m)(2+m)z&=&1-m.\\
\end{array}\right. $$
\begin{enumerate}
\item[-]Si $m=1$, $(\mathtt{S}) \Leftrightarrow  x+y+z=1 $ admet une infinité de solutions.
\item[-]Si $m=-2$, $(\mathtt{S})$ n'admet pas de solution.
\item[-]Si $m\neq 1$ et $m\neq -2$,  $(\mathtt{S})$ admet une unique solution.
\end{enumerate}
\end{explanations}
\end{question}
}
